\documentclass{article}
\usepackage[T1]{fontenc}
\usepackage{lmodern}
\usepackage{graphicx}
\usepackage{amsmath,amssymb}
\usepackage[a4paper,margin=2cm]{geometry}
\usepackage[absolute,overlay]{textpos}
\setlength{\TPHorizModule}{1mm}
\setlength{\TPVertModule}{1mm}
\usepackage[indonesian]{babel}
\usepackage{hyperref}
\setlength{\emergencystretch}{2em}
% unit: Halaman 36

\begin{document}

% === Halaman 36 (python/native) ===

36

2. Kode Program RSA Python (2)

from Crypto.PublicKey import RSA

from Crypto.Cipher import PKCS1\_OAEP

def PembangkitanKunciRSA():

    \# Buat kunci RSA

    key = RSA.generate(1024)

    kunci\_publik = key.publickey()


    \# Tampilkan informasi kunci

    print("Modulus (n):", key.n)

    print("Kunci publik (e):", key.e)

    print("Kunci privat (d):", key.d)

    print("Bilangan prima p:", key.p)

    print("Bilangan prima q:", key.q)

    \#Simpan kunci publik dan kunci privat ke dalam dua file kunci berbeda dengan format PEM

    with open("kunci\_publik.pem", "wb") as file:

         file.write(kunci\_publik.exportKey('PEM'))

         file.close()

    with open("kunci\_privat.pem", "wb") as file:

         file.write(key.exportKey('PEM','passwordku'))  \# Password bisa diubah

         file.close()

\end{document}
